\chapter{Căng Tin Và Giờ Ra Chơi}
\chapterintro{Không chỉ lớp học mới có chuyện để khịa}{Giờ ra chơi, đồ ăn, xếp hàng và những pha chạy chuông cũng là một phần rất thật của đời học sinh.}

\begin{lucbatbox}{Lục bát: Chuông ra chơi là tỉnh}
Trong giờ em mắt lim dim
Chuông vừa ra chơi bật dậy sống liền
Chân đi tốc độ thần tiên
Hướng về căng-tin nhẹ tênh cánh diều
Nếu mà học cũng bấy nhiêu
Thầy cô chắc đã thương yêu lắm rồi
\end{lucbatbox}

\begin{tutuyetbox}{Thất ngôn tứ tuyệt: Xếp hàng văn minh}
Miệng ai cũng bảo xếp hàng thôi
Chân thì nhích nhẹ tới từng người
Người sau bỗng thấy mình thành trước
Văn minh đôi lúc rẽ đường vui
\end{tutuyetbox}

\begin{ghichu}
Căng-tin là nơi bộc lộ rõ nhất tốc độ, chiến thuật, và sức mạnh sinh tồn của học sinh sau hai tiết liền.
\end{ghichu}

\ngancach

\begin{lucbatbox}{Lục bát: Mua theo phong trào}
Bạn kia mua món đang hot
Bạn này chưa đói cũng gật đầu theo
Một hồi bàn học xôn xao
Mùi cay mùi ngọt đua nhau kín bàn
Giờ vào còn tiếc dở dang
Miệng còn nhai nốt, mắt ngoan nhìn cô
\end{lucbatbox}

\begin{tutuyetbox}{Thất ngôn tứ tuyệt: Chạy chuông vào lớp}
Vừa nghe chuông đã cuống chân lên
Túi bánh trên tay mắt ngó ngang
Đời học trò vừa thơ vừa gấp
Nên lắm phen chạy mướt mồ hôi
\end{tutuyetbox}